% ═══════════════════════════════════════════════════════════════
% ML-01c: Artikel + Alphabet (MUSTERLÖSUNG)
% Spanisch Klasse 7 - Sequenz 1: Empezamos
% ═══════════════════════════════════════════════════════════════
% Identisches Layout wie AB-01c, Lösungen in ROT
% ═══════════════════════════════════════════════════════════════

\documentclass[11pt, a4paper]{article}

% ═══════════════════════════════════════════════════════════════
% SW-MODUS TOGGLE
% ═══════════════════════════════════════════════════════════════
\newif\ifswmodus
\swmodusfalse

% ═══════════════════════════════════════════════════════════════
% PAKETE
% ═══════════════════════════════════════════════════════════════

\usepackage[utf8]{inputenc}
\usepackage[T1]{fontenc}
\usepackage[ngerman]{babel}
\usepackage{helvet}
\renewcommand{\familydefault}{\sfdefault}

\usepackage[margin=2cm]{geometry}
\usepackage{enumitem}
\usepackage{tabularx}
\usepackage{booktabs}
\usepackage{multirow}

\usepackage{xcolor}
\usepackage{framed}
\usepackage{tcolorbox}
\tcbuselibrary{skins}

\usepackage{fancyhdr}
\usepackage{lastpage}
\usepackage{needspace}

% ═══════════════════════════════════════════════════════════════
% FARBEN
% ═══════════════════════════════════════════════════════════════

\definecolor{spanisch-color}{HTML}{c41e3a}
\definecolor{grau-color}{HTML}{888888}
\definecolor{hellgrau-color}{HTML}{f5f5f5}
\definecolor{orange-color}{HTML}{b35c00}
\definecolor{loesung-color}{HTML}{c62828}

\definecolor{sw-dark}{HTML}{000000}
\definecolor{sw-gray}{HTML}{666666}
\definecolor{sw-white}{HTML}{ffffff}

\ifswmodus
    \colorlet{spanisch}{sw-dark}
    \colorlet{grau}{sw-gray}
    \colorlet{hellgrau}{sw-white}
    \colorlet{orange}{sw-dark}
    \colorlet{loesung}{sw-dark}
\else
    \colorlet{spanisch}{spanisch-color}
    \colorlet{grau}{grau-color}
    \colorlet{hellgrau}{hellgrau-color}
    \colorlet{orange}{orange-color}
    \colorlet{loesung}{loesung-color}
\fi

\newcommand{\fachfarbe}{spanisch}

% ═══════════════════════════════════════════════════════════════
% VARIABLEN
% ═══════════════════════════════════════════════════════════════

\newcommand{\thema}{Artikel + Alphabet}
\newcommand{\sequenz}{01}
\newcommand{\block}{c}
\newcommand{\fach}{Spanisch}
\newcommand{\klasse}{7}
\newcommand{\maxpunkte}{20}
\newcommand{\datum}{\today}

% ═══════════════════════════════════════════════════════════════
% KOPF- UND FUSSZEILE
% ═══════════════════════════════════════════════════════════════

\pagestyle{fancy}
\fancyhf{}
\renewcommand{\headrulewidth}{0.4pt}
\renewcommand{\footrulewidth}{0.4pt}

\lhead{\fach{} -- Klasse \klasse}
\chead{\textcolor{loesung}{\textbf{ML-\sequenz\block{} (MUSTERLÖSUNG)}}}
\rhead{\datum}

\lfoot{}
\cfoot{}
\rfoot{Seite \thepage\ von \pageref{LastPage}}

% ═══════════════════════════════════════════════════════════════
% BEFEHLE
% ═══════════════════════════════════════════════════════════════

\newcommand{\punktebox}[1]{\hfill\fbox{\small \_/#1}}
\newcommand{\luecke}[1]{\rule{#1}{0.4pt}}
\newcommand{\lsg}[1]{\textcolor{loesung}{\textbf{#1}}}

\newtcolorbox{merkkasten}[1][]{
    colback=hellgrau,
    colframe=\fachfarbe,
    colbacktitle=white,
    coltitle=\fachfarbe,
    boxrule=1pt,
    arc=0pt,
    left=8pt, right=8pt, top=8pt, bottom=8pt,
    fonttitle=\bfseries,
    attach boxed title to top left={yshift=-2mm, xshift=5mm},
    boxed title style={boxrule=0pt, colback=white},
    title=#1
}

\newtcolorbox{vokabelkasten}{
    colback=white,
    colframe=\fachfarbe,
    boxrule=0.5pt,
    arc=0pt,
    left=5pt, right=5pt, top=5pt, bottom=5pt
}

\newtcolorbox{hinweis}{
    colback=white,
    colframe=orange,
    colbacktitle=white,
    coltitle=orange,
    boxrule=1pt,
    arc=0pt,
    left=5pt, right=5pt, top=5pt, bottom=5pt,
    fonttitle=\bfseries,
    attach boxed title to top left={yshift=-2mm, xshift=5mm},
    boxed title style={boxrule=0pt, colback=white},
    title=Hinweis
}

\newenvironment{aufgabe}[1]{%
    \needspace{5\baselineskip}%
    \nopagebreak[4]%
    \vspace{10pt}
    \noindent\textbf{\textcolor{\fachfarbe}{Aufgabe #1}}
    \nopagebreak[4]%
    \vspace{5pt}
    \nopagebreak[4]%
    \begin{list}{}{\leftmargin=0pt}
    \item[]
}{%
    \end{list}
}

\newcommand{\es}[1]{\textcolor{\fachfarbe}{\textbf{#1}}}

% ═══════════════════════════════════════════════════════════════
% DOKUMENT
% ═══════════════════════════════════════════════════════════════

\begin{document}

% ═══════════════════════════════════════════════════════════════
% KOPFBEREICH
% ═══════════════════════════════════════════════════════════════

\begin{center}
    {\Large\bfseries\textcolor{\fachfarbe}{Arbeitsblatt: \thema}}\\[0.3em]
    {\large\textcolor{loesung}{\textbf{--- MUSTERLÖSUNG ---}}}
\end{center}

\vspace{5pt}

\noindent
\begin{tabularx}{\textwidth}{|X|X|X|}
\hline
\textbf{Name:} \lsg{Musterschüler/in} &
\textbf{Klasse:} \klasse &
\textbf{Datum:} \datum \\
\hline
\end{tabularx}

\vspace{15pt}

% ═══════════════════════════════════════════════════════════════
% MERKKASTEN: Die bestimmten Artikel
% ═══════════════════════════════════════════════════════════════

\begin{merkkasten}[Merkkasten: Die bestimmten Artikel]
Im Spanischen gibt es \textbf{vier bestimmte Artikel}:

\vspace{0.5em}

\begin{center}
\begin{tabular}{|l|c|c|}
\hline
& \textbf{Singular} & \textbf{Plural} \\
\hline
\textbf{maskulin} & \es{el} \lsg{(der)} & \es{los} \lsg{(die)} \\
\hline
\textbf{feminin} & \es{la} \lsg{(die)} & \es{las} \lsg{(die)} \\
\hline
\end{tabular}
\end{center}

\vspace{0.5em}

\textbf{Genus-Regeln:}
\begin{itemize}[leftmargin=2em, itemsep=0pt]
    \item Nomen auf \textbf{-o} $\rightarrow$ meist maskulin (el libr\textbf{o})
    \item Nomen auf \textbf{-a} $\rightarrow$ meist feminin (la mes\textbf{a})
\end{itemize}
\end{merkkasten}

\vspace{10pt}

% ═══════════════════════════════════════════════════════════════
% AUFGABE 1: Zuordnung Nomen $\rightarrow$ Artikel
% ═══════════════════════════════════════════════════════════════

\begin{aufgabe}{1 -- Ordne die Nomen zu}
Schreibe die Wörter in die richtige Spalte. \punktebox{4}
\end{aufgabe}

\begin{vokabelkasten}
libro | mesa | cuaderno | silla | bolígrafo | pizarra | lápiz | ventana
\end{vokabelkasten}

\vspace{0.5em}

\begin{center}
\begin{tabular}{|p{6cm}|p{6cm}|}
\hline
\textbf{\es{el} (maskulin)} & \textbf{\es{la} (feminin)} \\
\hline
\lsg{libro, cuaderno,} & \lsg{mesa, silla,} \\
\lsg{bolígrafo, lápiz} & \lsg{pizarra, ventana} \\[1em]
\hline
\end{tabular}
\end{center}

\vspace{10pt}

% ═══════════════════════════════════════════════════════════════
% AUFGABE 2: Artikel einsetzen
% ═══════════════════════════════════════════════════════════════

\begin{aufgabe}{2 -- Setze den richtigen Artikel ein}
Wähle: \es{el}, \es{la}, \es{los} oder \es{las}. \punktebox{4}
\end{aufgabe}

\noindent
a) \lsg{el} chico \hfill
b) \lsg{la} profesora \\[0.8em]
c) \lsg{los} libros \hfill
d) \lsg{las} chicas \\

\vspace{15pt}

% ═══════════════════════════════════════════════════════════════
% MERKKASTEN: Das spanische Alphabet
% ═══════════════════════════════════════════════════════════════

\begin{merkkasten}[Merkkasten: Das spanische Alphabet]
\textbf{Besondere Buchstaben:}

\vspace{0.3em}

\begin{tabular}{ll}
\es{h} (hache) & stumm! (hola = [óla]) \\
\es{j} (jota) & wie deutsches CH (jamón) \\
\es{ñ} (eñe) & wie NJ (España) \\
\es{ll} (elle) & wie J (calle) \\
\es{rr} (erre) & gerolltes R (perro) \\
\end{tabular}
\end{merkkasten}

\vspace{10pt}

% ═══════════════════════════════════════════════════════════════
% AUFGABE 3: Buchstabieren
% ═══════════════════════════════════════════════════════════════

\begin{aufgabe}{3 -- Buchstabiere die Wörter}
Schreibe die Buchstaben auf Spanisch. \punktebox{6}
\end{aufgabe}

\noindent
a) \es{HOLA}: \lsg{hache - o - ele - a}

\vspace{0.5em}
\textit{Beispiel: hache - o - ele - a}

\vspace{1em}

\noindent
b) \es{GRACIAS}: \lsg{ge - erre - a - ce - i - a - ese}

\vspace{1.5em}

\noindent
c) \textbf{Dein Name}: \lsg{(individuell)}

\vspace{15pt}

% ═══════════════════════════════════════════════════════════════
% AUFGABE 4: Dialog mit Buchstabieren
% ═══════════════════════════════════════════════════════════════

\begin{aufgabe}{4 -- Vervollständige den Dialog}
A fragt nach dem Namen und wie man ihn schreibt. \punktebox{6}
\end{aufgabe}

\begin{hinweis}
Verwende: \es{¿Cómo te llamas?} und \es{¿Cómo se escribe?}
\end{hinweis}

\vspace{0.5em}

\noindent
\textbf{A:} \lsg{¿Cómo te llamas?}

\vspace{0.8em}

\noindent
\textbf{B:} Me llamo \lsg{María}.

\vspace{0.8em}

\noindent
\textbf{A:} \lsg{¿Cómo se escribe?}

\vspace{0.8em}

\noindent
\textbf{B:} \lsg{Eme - a - erre - i con acento - a}

\textit{(Buchstabiere deinen Namen!)}

% ═══════════════════════════════════════════════════════════════
% PUNKTEÜBERSICHT
% ═══════════════════════════════════════════════════════════════

\vspace{20pt}
\hrule
\vspace{10pt}

\begin{flushright}
\textbf{Punkte:} \lsg{\maxpunkte} / \maxpunkte
\end{flushright}

\end{document}
